\documentclass[11pt]{article}
\usepackage[margin=1in]{geometry}
\usepackage{hyperref}
\usepackage{listings}
\usepackage{xcolor}
\usepackage{booktabs}
\usepackage{enumitem}
\usepackage{parskip}
\usepackage{titlesec}

\hypersetup{
    colorlinks=true,
    linkcolor=blue,
    urlcolor=blue,
}

\definecolor{codebg}{rgb}{0.95,0.95,0.95}
\definecolor{codeframe}{rgb}{0.7,0.7,0.7}
\definecolor{kwcolor}{rgb}{0.0,0.3,0.7}
\definecolor{commentcolor}{rgb}{0.4,0.4,0.4}
\definecolor{stringcolor}{rgb}{0.1,0.5,0.1}

\lstset{
    backgroundcolor=\color{codebg},
    frame=single,
    rulecolor=\color{codeframe},
    basicstyle=\ttfamily\small,
    keywordstyle=\color{kwcolor}\bfseries,
    commentstyle=\color{commentcolor}\itshape,
    stringstyle=\color{stringcolor},
    breaklines=true,
    breakatwhitespace=true,
    showstringspaces=false,
    tabsize=4,
    xleftmargin=6pt,
    xrightmargin=6pt,
}

\title{\textbf{OSG/OSPool Quickstart Guide}\\
\large A100 GPU Job Submission for Fermionic Adaptive Circuit Simulations}
\author{Asadullah Bhuiyan \\ \small\texttt{ab2398@cornell.edu}}
\date{April 2026}

\begin{document}
\maketitle
\tableofcontents
\newpage

% ============================================================
\section{Account Setup and Login}
% ============================================================

After your OSG account is approved (within one business day of request), you
will be assigned to one of the UW access points. SSH in with:

\begin{lstlisting}[language=bash]
ssh asad.bhuiyan@ap40.uw.osg-htc.org
# alternates: ap41.uw.osg-htc.org
\end{lstlisting}

If SSH fails, upload your public key via the OSG portal web interface first,
then retry. To generate a key if you do not have one:

\begin{lstlisting}[language=bash]
ssh-keygen -t ed25519 -C "ab2398@cornell.edu"
# public key lives at ~/.ssh/id_ed25519.pub -- paste into portal
\end{lstlisting}

Check which project you belong to (needed for resource accounting):
\begin{lstlisting}[language=bash]
grep $USER /etc/condor/UserToProjectMap.txt
\end{lstlisting}

% ============================================================
\section{HTCondor Basics}
% ============================================================

OSPool uses HTCondor for job scheduling. The core workflow is:

\begin{enumerate}[leftmargin=*]
    \item Write an \textbf{executable} script (\texttt{run.sh} or \texttt{run.py}).
    \item Write a \textbf{submit file} (\texttt{*.submit}) declaring resources and
          pointing to the executable.
    \item Run \texttt{condor\_submit job.submit}.
    \item Monitor with \texttt{condor\_q} or \texttt{condor\_watch\_q}.
\end{enumerate}

\subsection{Minimal Submit File Structure}

\begin{lstlisting}
universe = container
container_image = /cvmfs/singularity.opensciencegrid.org/...

executable = run.sh
arguments  = $(Process)

log    = logs/job_$(Cluster)_$(Process).log
error  = logs/job_$(Cluster)_$(Process).err
output = logs/job_$(Cluster)_$(Process).out

transfer_input_files = run.sh

+JobDurationCategory = "Medium"

request_gpus   = 1
request_cpus   = 4
request_memory = 16 GB
request_disk   = 10 GB

queue 1
\end{lstlisting}

\textbf{Important:} \texttt{+JobDurationCategory} is \emph{required} on all
OSPool access points. Missing it causes the job to be held immediately.

\subsection{Job Duration Categories}

\begin{center}
\begin{tabular}{ll}
\toprule
Category & Wall time \\
\midrule
\texttt{"Short"}  & $< 1$ hour \\
\texttt{"Medium"} & $1$--$10$ hours \\
\texttt{"Long"}   & $10$--$20$ hours \\
\bottomrule
\end{tabular}
\end{center}

Jobs exceeding 20 hours require HTCondor checkpointing or they will be evicted.

\subsection{Monitoring Commands}

\begin{lstlisting}[language=bash]
condor_q                          # queue summary
condor_watch_q                    # live update (Ctrl+C to exit)
condor_q -hold                    # see why held jobs are held
condor_q -format "%d\n" JobStatus | sort | uniq -c
# status codes: 1=Idle, 2=Running, 4=Complete, 5=Held

condor_rm <cluster-id>            # remove entire cluster
condor_rm <cluster-id>.<process>  # remove single job
condor_rm $USER                   # remove all your jobs
\end{lstlisting}

% ============================================================
\section{Targeting A100 GPUs}
% ============================================================

A100 GPUs have CUDA Compute Capability 8.0. Use these HTCondor ClassAd
attributes to target them:

\begin{lstlisting}
# Pin exactly to A100 (CC = 8.0):
gpus_minimum_capability = 8.0
gpus_maximum_capability = 8.0

# Require at least 40 GB GPU RAM (A100-40GB variant):
gpus_minimum_memory = 40000

# For the 80 GB variant:
gpus_minimum_memory = 80000
\end{lstlisting}

\subsection{GPU Availability vs.\ Specificity Trade-off}

\begin{center}
\begin{tabular}{lll}
\toprule
Constraint & Hits & Notes \\
\midrule
\texttt{capability = 8.0..8.0} & A100 only & Most restrictive, may queue longer \\
\texttt{capability $\geq$ 8.0} & A100 + A10 (22 GB) & Good throughput \\
\texttt{capability $\geq$ 7.0} & + V100 (16 GB) & Widest pool \\
\bottomrule
\end{tabular}
\end{center}

For \texttt{classA\_U1FGTN\_gpu} runs with \texttt{complex128} covariance
matrices at $N_x{=}20$, $N_y{\leq}120$, the A100-40GB is recommended.
The V100 (16 GB) may be insufficient for large $N_y$.

% ============================================================
\section{Container Setup for Your Environment}
% ============================================================

Your code requires packages not present in any pre-built OSPool container
(\texttt{pfapack}, \texttt{polars}, \texttt{h5py}, \texttt{torch 2.8},
\texttt{scipy}, etc.). You must build a custom container.

\subsection{Pre-Built Containers (Reference)}

\begin{center}
\begin{tabular}{ll}
\toprule
CVMFS Path & Contents \\
\midrule
\texttt{.../htc/pytorch:2.3.1-cuda11.8} & PyTorch 2.3.1, CUDA 11.8 \\
\texttt{.../htc/tensorflow:2.15}         & TensorFlow 2.15 \\
\texttt{.../htc/rocky:9-cuda-12.6.0}     & Rocky 9 + CUDA 12.6 (bare) \\
\texttt{.../htc/rocky:8-cuda-11.0.3}     & Rocky 8 + CUDA 11.0 (bare) \\
\bottomrule
\end{tabular}
\end{center}

Full CVMFS prefix: \texttt{/cvmfs/singularity.opensciencegrid.org/}

\subsection{Building a Custom Container (Local Machine)}

Docker is \emph{not} available on OSPool access points. Build locally, push
to Docker Hub, then pull on the access point.

\textbf{Step 1 --- Dockerfile} (save as \texttt{Dockerfile}):

\begin{lstlisting}[language=bash]
FROM hub.opensciencegrid.org/htc/rocky:9-cuda-12.6.0

# system deps
RUN dnf install -y python3 python3-pip && dnf clean all

# install your exact pinned environment
RUN pip install --no-cache-dir \
    torch==2.8.0 \
    numpy \
    scipy \
    matplotlib \
    h5py \
    pfapack \
    tqdm \
    pandas \
    polars \
    seaborn \
    fastparquet \
    pyarrow \
    opt-einsum \
    joblib \
    cloudpickle \
    networkx \
    sympy
\end{lstlisting}

\textbf{Step 2 --- Build and push:}

\begin{lstlisting}[language=bash]
docker build -t <dockerhub-username>/fermionic-env:latest .
docker push <dockerhub-username>/fermionic-env:latest
\end{lstlisting}

\textbf{Step 3 --- Convert to \texttt{.sif} on the access point:}

\begin{lstlisting}[language=bash]
# On ap40 (or ap41):
mkdir -p $HOME/tmp
export TMPDIR=$HOME/tmp
export APPTAINER_TMPDIR=$HOME/tmp
export APPTAINER_CACHEDIR=$HOME/tmp

apptainer build fermionic-env.sif \
    docker://<dockerhub-username>/fermionic-env:latest
\end{lstlisting}

For images $>1$ GB, stage through OSDF rather than \texttt{transfer\_input\_files}.
For quick iteration, reference the Docker Hub image directly in the submit file:

\begin{lstlisting}
container_image = docker://<dockerhub-username>/fermionic-env:latest
\end{lstlisting}

% ============================================================
\section{Submitting 100 A100 Jobs}
% ============================================================

\subsection{Submit File: \texttt{batch\_gpu.submit}}

\begin{lstlisting}
universe = container
container_image = fermionic-env.sif

executable = run.sh
arguments  = $(Process)

transfer_input_files = run.sh, fermionic-env.sif

log    = logs/job_$(Cluster)_$(Process).log
error  = logs/job_$(Cluster)_$(Process).err
output = logs/job_$(Cluster)_$(Process).out

+JobDurationCategory = "Medium"

# Target A100 (Compute Capability 8.0, 40+ GB VRAM)
gpus_minimum_capability = 8.0
gpus_maximum_capability = 8.0
gpus_minimum_memory     = 40000

request_gpus   = 1
request_cpus   = 4
request_memory = 16 GB
request_disk   = 10 GB

queue 100
\end{lstlisting}

\textbf{Note:} If you use a CVMFS path or a \texttt{docker://} URI for
\texttt{container\_image}, remove it from \texttt{transfer\_input\_files}.
Only include local \texttt{.sif} files there.

\subsection{Executable: \texttt{run.sh}}

\texttt{\$(Process)} runs 0--99. Use it to shard your parameter sweep or
sample range:

\begin{lstlisting}[language=bash]
#!/bin/bash
JOB_ID=$1

echo "Job ${JOB_ID} starting on $(hostname)"
echo "GPU: $(nvidia-smi --query-gpu=name --format=csv,noheader)"

python3 run_job.py --job-id ${JOB_ID}
\end{lstlisting}

\begin{lstlisting}[language=bash]
chmod +x run.sh
mkdir -p logs
condor_submit batch_gpu.submit
\end{lstlisting}

\subsection{Important Batch Submission Notes}

\begin{itemize}[leftmargin=*]
    \item 100 jobs is well under the per-user 10,000-job queue limit.
    \item Jobs start as slots open across the pool, not simultaneously.
    \item Always use both \texttt{\$(Cluster)} and \texttt{\$(Process)} in
          log/output paths or files will overwrite each other.
    \item To cap how many run at once (e.g.\ to avoid storage contention):
          \texttt{max\_idle = 20} in the submit file.
\end{itemize}

% ============================================================
\section{Mapping Your Workflows to OSPool Jobs}
% ============================================================

\subsection{Pure-State Markov Snapshot Runs (\texttt{colab\_small\_system\_testing})}

These notebooks call \texttt{classA\_U1FGTN\_gpu.run\_markov\_circuit()} for
configurations like $N_x{=}16$, $N_y{\in}\{30,40\}$, \texttt{nshell}$\in\{1,2\}$,
10 samples, 50 cycles. Each configuration is independent---natural to
parallelize as separate HTCondor jobs.

A parameter sweep over \texttt{(Ny, nshell)} might look like:

\begin{lstlisting}[language=python]
# generate_params.py  -- run locally to create params.json
import json
params = []
for Ny in [30, 40]:
    for nshell in [1, 2]:
        params.append({"Ny": Ny, "nshell": nshell})
with open("params.json", "w") as f:
    json.dump(params, f)
\end{lstlisting}

\begin{lstlisting}[language=bash]
#!/bin/bash
# run.sh -- receives job index as $1
JOB_ID=$1
python3 -c "
import json, sys
params = json.load(open('params.json'))
p = params[int('${JOB_ID}')]
print(p)
# call classA_U1FGTN_gpu with p['Ny'], p['nshell'], etc.
"
\end{lstlisting}

\subsection{Large $N{=}20$ Entanglement Scaling (\texttt{colab\_large\_entanglement\_scaling\_N20})}

These notebooks run $N_x{=}20$, \texttt{samples=100}, $N_y\in\{30,40,50,60,80,100,120\}$,
\texttt{nshell=1}, 50 cycles. Each $N_y$ is independent. With 100 samples per
$N_y$, you can shard samples across jobs:

\begin{lstlisting}
# Submit 7 jobs, one per Ny value:
queue 7

# Or shard 100 samples as 10 jobs x 10 samples per Ny:
queue 70   # 7 Ny values * 10 jobs each
\end{lstlisting}

The GPU driver uses \texttt{complex128} by default. For $N_x{=}20$, $N_y{=}120$,
the covariance matrix is $4 N_x N_y / 2 = 4800$ dimensional---comfortably
within A100-40GB RAM.

\subsection{Strip Entanglement Post-Processing}

\texttt{strip\_entropy\_contours\_gpu.py} and
\texttt{strip\_entropy\_streaming\_gpu.py} perform observable post-processing
on saved covariance snapshots. These are also GPU-bound (matrix diagonalization
via CUDA) and can be submitted as separate jobs after the snapshot runs
complete, using DAGMan for dependency tracking:

\begin{lstlisting}
# simple.dag
JOB  SNAPSHOTS  snapshot_jobs.submit
JOB  STRIPS     strip_jobs.submit
PARENT SNAPSHOTS CHILD STRIPS
\end{lstlisting}

\begin{lstlisting}[language=bash]
condor_submit_dag simple.dag
\end{lstlisting}

% ============================================================
\section{Data Transfer}
% ============================================================

\subsection{Small Files ($<$100 MB): \texttt{transfer\_input\_files}}

\begin{lstlisting}
transfer_input_files = run.sh, params.json, src/classA_U1FGTN_gpu.py
transfer_output_files = output_$(Process).npz
\end{lstlisting}

\subsection{Larger Files: OSDF}

For \texttt{.sif} containers, large \texttt{.npy}/\texttt{.npz} snapshots, or
output data $>$100 MB, use the Open Science Data Federation (OSDF):

\begin{lstlisting}[language=bash]
# Upload to OSDF staging area from access point:
osdf object put fermionic-env.sif \
    /osdf/users/asad.bhuiyan/fermionic-env.sif

# Reference in submit file:
container_image = osdf:///osdf/users/asad.bhuiyan/fermionic-env.sif
\end{lstlisting}

\subsection{Transferring Results Back}

\begin{lstlisting}[language=bash]
# From your local machine or access point after jobs finish:
scp asad.bhuiyan@ap40.uw.osg-htc.org:~/results/ ./results/ -r
\end{lstlisting}

% ============================================================
\section{Common Pitfalls and Fixes}
% ============================================================

\begin{center}
\begin{tabular}{p{0.35\textwidth}p{0.55\textwidth}}
\toprule
Problem & Fix \\
\midrule
Job held immediately & Missing \texttt{+JobDurationCategory}. Add it. \\
Job runs on CPU node & Missing \texttt{request\_gpus = 1}. \\
No A100s available & Remove \texttt{gpus\_maximum\_capability} to open up A10s too, or lower \texttt{gpus\_minimum\_capability} to 7.0 for V100s. \\
Output files overwrite each other & Use \texttt{\$(Cluster)\_\$(Process)} in all output paths. \\
Job evicted after 20 hours & Implement HTCondor checkpointing or split into shorter segments. \\
\texttt{.sif} too large to transfer & Stage via OSDF instead of \texttt{transfer\_input\_files}. \\
\texttt{torch.cuda.is\_available()} returns \texttt{False} & The container must be GPU-capable and \texttt{request\_gpus $\geq$ 1}. \\
\texttt{pfapack} not found & Must build custom container; not in any pre-built OSPool image. \\
\bottomrule
\end{tabular}
\end{center}

% ============================================================
\section{Full Working Example}
% ============================================================

This is a self-contained example for running 100 A100 jobs with your
\texttt{fermionic\_env} container.

\textbf{\texttt{run.sh}:}
\begin{lstlisting}[language=bash]
#!/bin/bash
set -euo pipefail
JOB_ID=$1

echo "=== Job ${JOB_ID} on $(hostname) ==="
nvidia-smi --query-gpu=name,memory.total --format=csv,noheader

python3 - <<'EOF'
import sys, os, torch
import numpy as np

job_id = int(os.environ.get("JOB_ID", sys.argv[1] if len(sys.argv) > 1 else "0"))
device = torch.device("cuda" if torch.cuda.is_available() else "cpu")
print(f"Job {job_id} | device: {device} | GPU: {torch.cuda.get_device_name(0)}")

# --- your actual work here ---
# e.g.: import classA_U1FGTN_gpu; run_markov_circuit(...)
# save outputs to output_{job_id}.npz

np.savez_compressed(f"output_{job_id}.npz", job_id=job_id)
print(f"Job {job_id} done.")
EOF
\end{lstlisting}

\textbf{\texttt{batch\_gpu.submit}:}
\begin{lstlisting}
universe = container
container_image = docker://<dockerhub-username>/fermionic-env:latest

executable = run.sh
arguments  = $(Process)

environment = "JOB_ID=$(Process)"

transfer_input_files = run.sh

log    = logs/job_$(Cluster)_$(Process).log
error  = logs/job_$(Cluster)_$(Process).err
output = logs/job_$(Cluster)_$(Process).out

+JobDurationCategory = "Medium"

gpus_minimum_capability = 8.0
gpus_maximum_capability = 8.0
gpus_minimum_memory     = 40000

request_gpus   = 1
request_cpus   = 4
request_memory = 16 GB
request_disk   = 10 GB

queue 100
\end{lstlisting}

\textbf{Submit:}
\begin{lstlisting}[language=bash]
chmod +x run.sh
mkdir -p logs
condor_submit batch_gpu.submit
condor_watch_q
\end{lstlisting}

% ============================================================
\section{Quick Reference Card}
% ============================================================

\begin{lstlisting}[language=bash]
# Login
ssh asad.bhuiyan@ap40.uw.osg-htc.org

# Submit
condor_submit job.submit

# Monitor
condor_q
condor_watch_q
condor_q -hold               # why are jobs held?

# Cancel
condor_rm <cluster-id>
condor_rm $USER              # cancel everything

# GPU ClassAds (A100)
gpus_minimum_capability = 8.0
gpus_maximum_capability = 8.0
gpus_minimum_memory     = 40000

# Required tag
+JobDurationCategory = "Medium"   # Short / Medium / Long

# Pre-built GPU containers (CVMFS, no transfer needed)
# /cvmfs/singularity.opensciencegrid.org/htc/pytorch:2.3.1-cuda11.8
# /cvmfs/singularity.opensciencegrid.org/htc/rocky:9-cuda-12.6.0

# Build custom .sif
export TMPDIR=$HOME/tmp APPTAINER_TMPDIR=$HOME/tmp APPTAINER_CACHEDIR=$HOME/tmp
apptainer build myenv.sif docker://<user>/<image>:tag

# Help
support@osg-htc.org
\end{lstlisting}

\end{document}
